%% 1-4 たし算は「あわせる」「ふえる」（10ページ、p.46〜55）
%% 第1巻のきまり：15 pt・行送り28 pt、分かち書き（文節の間に半角の空白）、漢字はすべてルビ、式は全角の＋＝で書く

\bkchapter{1}{4}{たし\R{算}{ざん}は\\「あわせる」「ふえる」}{\bklevel{1}}{たしざんは「あわせる」「ふえる」}
\bkmichishirube{1-1,1-2}{1-5,1-7,1-12,2-3}

\begin{goal}
\begin{itemize}[label={\color{bkTheme}●}, leftmargin=1.4\zw, itemsep=0pt, topsep=0pt]
\item 「あわせる」「ふえる」を、 たし\R{算}{ざん}の しきに かける。
\item ものさしの うえで、 たし\R{算}{ざん}が できる。
\item 10を こえる たし\R{算}{ざん}が できる。
\end{itemize}
\end{goal}

\begin{hatena}
\begin{bkfigure}
\begin{tikzpicture}[x=1mm, y=1mm]
  % あ：おさらと かご
  \fill[white, rounded corners=3mm] (0,0) rectangle (59,44);
  \node[circle, fill=bkYellow, text=white, font=\sffamily\bfseries, minimum size=6.5mm, inner sep=0pt] at (6,38) {A};
  \bkplate{(16,13)}{24}
  \foreach \x/\y in {9.5/16.5,16/18,22.5/16.5} {\bkapple{(\x,\y)}{3.6}{bkAppleRed}}
  \draw[bkStem, line width=1.2pt, fill=bkYellow!25, rounded corners=1.5mm] (34,6) -- (54,6) -- (56,20) -- (32,20) -- cycle;
  \draw[bkStem, line width=1.2pt] (35,20) .. controls (36,31) and (52,31) .. (53,20);
  \foreach \x in {40,48} {\bkapple{(\x,23.5)}{3.6}{bkAppleRed}}
  \node[anchor=north west, font=\small, text=bkInk, align=left] at (13,42) {おさらに 3こ、\\かごに 2こ};
  % い：はこに 2こ いれる
  \begin{scope}[shift={(63,0)}]
  \fill[white, rounded corners=3mm] (0,0) rectangle (59,44);
  \node[circle, fill=bkYellow, text=white, font=\sffamily\bfseries, minimum size=6.5mm, inner sep=0pt] at (6,38) {B};
  \draw[bkStem, line width=1.2pt, fill=bkYellow!30] (6,4) rectangle (36,17);
  \fill[bkStem!40] (6,17) -- (10,21) -- (40,21) -- (36,17) -- cycle;
  \draw[bkStem, line width=1.2pt] (6,17) -- (10,21) -- (40,21) -- (36,17) (36,4) -- (40,8) -- (40,21);
  \foreach \x in {13,22,31} {\bkapple{(\x,22)}{3.6}{bkAppleRed}}
  \foreach \x in {47,55} {\bkapple{(\x,23)}{3.6}{bkAppleRed}}
  \draw[bk arrow, bkGray, line width=1pt] (49,18) -- (42.5,13.5);
  \node[anchor=north west, font=\small, text=bkInk, align=left] at (13,42) {はこに 3こ。\\あとから 2こ いれる};
  \end{scope}
\end{tikzpicture}
\end{bkfigure}
\noindent A も B も、 リンゴは ぜんぶで 5こ。\\
でも、 A は 2つを あわせる。 B は あとから ふえる。\\
ちがう ことなのに、 おなじ たし\R{算}{ざん}で いいのかな？
\end{hatena}

\clearpage
%% ---------------- p.47 やってみよう ----------------
\begin{yattemiyou}[おはじきで たしかめよう]
\begin{enumerate}[label={\bkcircled[bkGreen]{\arabic*}}, leftmargin=2.2\zw, itemsep=2pt, topsep=0pt]
\item \R{左}{ひだり}の おさらに おはじきを 3こ、 \R{右}{みぎ}の おさらに 2こ おこう。
\item 2つの おさらの おはじきを、 まんなかの おさらに あつめて かぞえよう。
\end{enumerate}
\begin{bkfigure}
\begin{tikzpicture}[x=1mm, y=1mm]
  \bkplate{(14,10)}{26}\bkplate{(104,10)}{26}\bkplate{(59,-6)}{32}
  \foreach \x in {7,14,21} {\bkohajiki{(\x,11)}{2.6}{bkRoute2}}
  \foreach \x in {100,108} {\bkohajiki{(\x,11)}{2.6}{bkRoute4}}
  \foreach \x in {48,55,62} {\bkohajiki{(\x,-5)}{2.6}{bkRoute2}}
  \foreach \x in {66,73} {\bkohajiki{(\x,-4)}{2.6}{bkRoute4}}
  \draw[bk arrow, bkGreen] (26,6) -- (41,-1);
  \draw[bk arrow, bkGreen] (92,6) -- (77,-1);
  \node[font=\sffamily\bfseries\small, text=bkGreen!70!black] at (14,20) {3こ};
  \node[font=\sffamily\bfseries\small, text=bkGreen!70!black] at (104,20) {2こ};
  \node[font=\sffamily\bfseries\small, text=bkGreen!70!black] at (59,-16) {あわせて なんこ？};
\end{tikzpicture}
\end{bkfigure}
\begin{enumerate}[label={\bkcircled[bkGreen]{\arabic*}}, leftmargin=2.2\zw, itemsep=2pt, topsep=0pt, start=3]
\item こんどは、 まんなかの おさらに 3こ おいてから、 1こずつ 2こ ふやして かぞえよう。
\end{enumerate}
\begin{bkfigure}
\begin{tikzpicture}[x=1mm, y=1mm]
  \bkplate{(40,0)}{40}
  \foreach \x in {28,35,42} {\bkohajiki{(\x,1)}{2.6}{bkRoute2}}
  \bkohajiki{(52,3)}{2.6}{bkRoute4}\bkohajiki{(60,15)}{2.6}{bkRoute4}
  \draw[bk arrow, bkGreen] (62,12.5) -- (55,5.5);
  \node[anchor=west, font=\sffamily\bfseries\small, text=bkGreen!70!black] at (66,15) {1こ、 また 1こ…};
\end{tikzpicture}
\end{bkfigure}
\noindent ②と ③で、 おはじきは なんこに なったかな？\par\vspace{3pt}
\noindent\hspace*{1\zw}②の とき\enspace\bkkakikomi\enspace こ\hspace{3\zw}③の とき\enspace\bkkakikomi\enspace こ
\end{yattemiyou}

\bkfukidashi{mina}{どっちも 5こに なったよ！ 「あわせる」と 「ふえる」は、 にて いるね。}

\clearpage
%% ---------------- p.48〜49 数学の道具 ----------------
\bksection{doguu}{\R{数学}{すう|がく}の \R{道具}{どう|ぐ}}

\bksubsection{1}{「あわせる」 たし\R{算}{ざん}}

\label{bk:1-4:awaseru}リンゴ 3こと 2こを あわせると 5こ。 しきで こう かきます。

\begin{bkfigure}
\begin{tikzpicture}[x=1mm, y=1mm]
  \foreach \x in {0,12,24} {\bkapple{(\x,0)}{4.6}{bkAppleRed}}
  \foreach \x in {48,60} {\bkapple{(\x,0)}{4.6}{bkAppleGreen}}
  \node[font=\fontsize{20}{20}\selectfont\bfseries, text=bkTheme] at (36,0) {＋};
  \node[font=\sffamily\bfseries\small, text=bkAppleRed!80!black] at (12,8.5) {3こ};
  \node[font=\sffamily\bfseries\small, text=bkAppleGreen!60!black] at (54,8.5) {2こ};
  \node[font=\fontsize{20}{20}\selectfont\bfseries, text=bkTheme] at (74,0) {→};
  \node[anchor=west, font=\sffamily\bfseries] at (80,0) {あわせて 5こ};
\end{tikzpicture}
\end{bkfigure}
\bkbigshiki{3 ＋ 2 ＝ 5}{3 たす 2 は 5}
\noindent\begin{minipage}[c]{0.55\linewidth}
「＋」は 「たす」と よみます。\bkidx{たしざん}{たし\R{算}{ざん}}\bkidx{たす}{＋（たす）}\bkidx{あわせる}{あわせる}\par
てんびんが つりあう ように、 「＝」の \R{左}{ひだり}と \R{右}{みぎ}は おなじ かずです。\bkidx{は}{＝（は）}\par
2つを いれかえて 2 ＋ 3 と かいても、 あわせると 5こ。 あわせる ときは、 どちらを さきに かいても おなじです。
\end{minipage}\hfill
\begin{minipage}[c]{0.43\linewidth}\centering
\begin{tikzpicture}[x=0.85mm, y=0.85mm]
  \bkbalance{(0,0)}{0}
  \foreach \x in {-29.4,-24,-18.6} {\bkapple{(\x,18.3)}{2.7}{bkAppleRed}}
  \foreach \x in {-26.7,-21.3} {\bkapple{(\x,23)}{2.7}{bkAppleGreen}}
  \foreach \x in {18.6,24,29.4} {\bkapple{(\x,18.3)}{2.7}{bkAppleRed}}
  \foreach \x in {21.3,26.7} {\bkapple{(\x,23)}{2.7}{bkAppleRed}}
  \node[anchor=base, font=\sffamily\bfseries\small] at (-24,-7) {3 ＋ 2};
  \node[anchor=base, font=\sffamily\bfseries\small, text=bkTheme] at (0,-7) {＝};
  \node[anchor=base, font=\sffamily\bfseries\small] at (24,-7) {5};
\end{tikzpicture}
\end{minipage}

\bksubsection{2}{「ふえる」 たし\R{算}{ざん}}

\label{bk:1-4:fueru}バスに 4\R{人}{にん} のって います。\bkidx{ふえる}{ふえる} 2\R{人}{にん} のって くると、 なん\R{人}{にん}？

\begin{bkfigure}
\begin{tikzpicture}[x=1mm, y=1mm]
  \bkbus{(0,0)}{4}{bkRoute6}
  \foreach \x in {92,100} {\bkperson{(\x,0)}{17}{bkRoute2}}
  \draw[bk arrow, bkRoute2] (104,22) -- (84,22);
  \node[font=\sffamily\bfseries\small, text=bkRoute2!80!black] at (95,27) {2\R{人}{にん} のる};
\end{tikzpicture}
\end{bkfigure}

\noindent 4\R{人}{にん}が 2\R{人}{にん} ふえて 6\R{人}{にん}。 これも たし\R{算}{ざん}です。
\bkbigshiki{4 ＋ 2 ＝ 6}{4 たす 2 は 6}

%% （第11回：ここで ページを かえない）
\bksubsection{3}{ものさしの うえを すすむ}

\label{bk:1-4:monosashi}3の ところから 2つ すすむと 5。 これも 3 ＋ 2 ＝ 5 です。

\begin{bkfigure}
\begin{tikzpicture}[x=1mm, y=1mm]
  \bkruler{(0,0)}{11}{10}
  \bkhop{(0,0)}{11}{3}{2}{bkTheme}
  \node[font=\sffamily\bfseries\small, text=bkTheme] at (44,14) {2つ すすむ};
\end{tikzpicture}
\end{bkfigure}
\vspace{3mm}
\begin{taisetsu}
「あわせる」「ふえる」「すすむ」は、 \\
どれも たし\R{算}{ざん}の しきに なります。
\end{taisetsu}

\noindent \bkhatenaref の A は 「あわせる」、 B は 「ふえる」。 どちらも 3 ＋ 2 ＝ 5 です。

\bkkigocards{\bkkigocard{＋}{たす}{あわせる、 ふえる、 すすむ ときの しるし}{plus}{1-4}}
{\bkkigocard{＝}{は}{\R{左}{ひだり}と \R{右}{みぎ}が おなじ かずの しるし}{equals}{1-4}}


%% （第11回：ここで ページを かえない）
%% ---------------- 第11回：いくつと いくつ（10を こえる たし算の じゅんび） ----------------
\bksubsection{4}{いくつと いくつ}

\label{bk:1-4:ikutsu}おはじき 5こを、 \R{線}{せん}で 2つに わけて みよう。

\begin{bkfigure}
\begin{tikzpicture}[x=1mm, y=1mm]
  \foreach \a [count=\k from 0] in {1,2,3,4} {
    \begin{scope}[shift={(\k*33,0)}]
      \fill[white, draw=bkLine, line width=0.8pt, rounded corners=2mm] (0,0) rectangle (30,20);
      \foreach \j in {1,...,5} {\ifnum\j>\a \def\cc{bkRoute4}\else\def\cc{bkRoute2}\fi
        \pgfmathsetmacro\xx{\j*5+(\j>\a ? 1.6 : -1.6)}\bkohajiki{(\xx-0.2,13)}{2}{\cc}}
      \draw[bkTheme, line width=1.2pt] ({\a*5+2.3},8.5) -- ({\a*5+2.3},17.5);
      \pgfmathtruncatemacro\b{5-\a}
      \node[font=\sffamily\bfseries\small] at (15,4.2) {\a と \b};
    \end{scope}}
\end{tikzpicture}
\end{bkfigure}

\noindent 5は、 1と 4、 2と 3、 3と 2、 4と 1に わけられます。 わけても、 あわせると もとの 5に もどります。 わける ばしょを かえても、 おはじきの かずは かわらないからです。

\noindent 10も おなじように わけられます。 10の わく（1-1）で たしかめよう。 あおと みどりを あわせると、 どれも わくが いっぱいの 10です。

\begin{bkfigure}
\begin{tikzpicture}[x=1mm, y=1mm]
  \foreach \a [count=\k from 0] in {1,...,9} {
    \pgfmathtruncatemacro\cx{mod(\k,3)}\pgfmathtruncatemacro\cy{div(\k,3)}
    \pgfmathtruncatemacro\b{10-\a}
    \begin{scope}[shift={(\cx*44,-\cy*17)}]
      \bktenframe{(0,0)}{3.6}{\a}{bkRoute2}{\b}{bkRoute4}
      \node[anchor=west, font=\sffamily\bfseries\small] at (20,3.6) {\a と \b};
    \end{scope}}
\end{tikzpicture}
\end{bkfigure}

\noindent この 「10に なる くみ」を おぼえて おくと、 つぎの 10を こえる たし\R{算}{ざん}が かんたんに なります。 たとえば 「8は あと 2で 10」です。\bkidx{じゅうになるくみ}{10に なる くみ}

\clearpage
\bksubsection{5}{10を こえる たし\R{算}{ざん}}

\label{bk:1-4:kuriagari}あめが 8こ。 5こ もらうと、 ぜんぶで なんこ？ 「10の まとまり」（1-2）を つくろう。

\noindent\begin{minipage}[c]{0.53\linewidth}
\begin{tikzpicture}[x=1mm, y=1mm]
  \bktenframe{(0,0)}{8}{8}{bkRoute2}{2}{bkRoute4}
  \foreach \x in {54,61,68} {\bkohajiki{(\x,4)}{3}{bkRoute4}}
  \draw[bk arrow, bkRoute4] (49.5,4) -- (41.5,4);
  \node[font=\sffamily\bfseries\footnotesize, text=bkRoute4!70!black] at (45.5,10) {2こ};
  \node[anchor=north, font=\sffamily\bfseries\footnotesize] at (20,-1.5) {10の まとまり};
  \node[anchor=north, font=\sffamily\bfseries\footnotesize] at (61,-1.5) {のこり 3こ};
\end{tikzpicture}
\end{minipage}\hfill
\begin{minipage}[c]{0.44\linewidth}
\begin{enumerate}[label={\bkcircled[bkTheme]{\arabic*}}, leftmargin=2\zw, itemsep=0pt, topsep=0pt]
\item 8は、 あと 2で 10。
\item 5を 2と 3に わける。
\item 8と 2で 10。
\item 10と 3で 13。
\end{enumerate}
\end{minipage}
\bkbigshiki{8 ＋ 5 ＝ 13}{8 たす 5 は 13}

\noindent どうして 5を 2と 3に わけるのかな？ 8を 10に するには、 あと 2 いります（10に なる くみ）。 だから、 5の なかから 2を 8に わたします。 5は 2と 3なので（いくつと いくつ）、 のこりは 3。 10と 3で 13（じゅうさん、 1-2）です。

\noindent ものさしの うえを すすんでも、 たしかめられます。 8の ところから 5つ すすむと……

\begin{bkfigure}
\begin{tikzpicture}[x=1mm, y=1mm]
  \bkruler{(0,0)}{8}{15}
  \bkhop{(0,0)}{8}{8}{2}{bkRoute2}
  \bkhop{(16,0)}{8}{8}{3}{bkRoute4}
  \fill[bkRoute2] (64,3.2) circle (1.6);
  \draw[bkTheme, line width=1.2pt] (104,-5.6) circle (3.4);
  \node[font=\sffamily\bfseries\footnotesize, text=bkRoute2!70!black] at (72,13) {2つ};
  \node[font=\sffamily\bfseries\footnotesize, text=bkRoute4!70!black] at (92,13) {3つ};
\end{tikzpicture}
\end{bkfigure}

\noindent 2つ すすむと 10。 そこから 3つ すすむと 13。 やっぱり 13に つきました。

\bksubsection{6}{0を たす たし\R{算}{ざん}}

\label{bk:1-4:zero}たまいれを 2かい しました。 1かいめは 3こ、 2かいめは 0こ。 あわせて なんこ？

\noindent\begin{minipage}[c]{0.5\linewidth}\centering
\begin{tikzpicture}[x=1mm, y=1mm]
  \bkplate{(14,4)}{24}\bkplate{(44,4)}{24}
  \foreach \x in {8,14,20} {\bkohajiki{(\x,6.5)}{2.5}{bkAppleRed}}
  \node[font=\sffamily\bfseries\small] at (14,-5) {1かいめ 3こ};
  \node[font=\sffamily\bfseries\small] at (44,-5) {2かいめ 0こ};
  \node[font=\sffamily\bfseries\small, text=bkTheme] at (29,4) {＋};
\end{tikzpicture}
\end{minipage}\hfill
\begin{minipage}[c]{0.47\linewidth}
0は、 なにも ない かず（1-3）。 なにも ふえないので、 かずは 3の ままです。
\end{minipage}
\bkbigshiki{3 ＋ 0 ＝ 3}{3 たす 0 は 3}
1かいめが 0こで、 2かいめが 3こでも、 0 ＋ 3 ＝ 3 です。 0を たしても、 0に たしても、 かずは かわりません。



%% （第11回：ここで ページを かえない）
%% ---------------- p.50〜51 例題 ----------------

\begin{reidai}[あわせる]
あかい ふうせんが 4こ、 あおい ふうせんが 3こ。 あわせて なんこ？

\noindent\begin{minipage}[c]{0.66\linewidth}
\kangaekata あわせるから、 たし\R{算}{ざん}。

\shiki 4 ＋ 3 ＝ 7 \kotaenext 7こ
\end{minipage}\hfill
\begin{minipage}[c]{0.32\linewidth}\centering
\begin{tikzpicture}[x=1mm, y=1mm]
  \foreach \x in {0,5,10,15} {\bkballoon{(\x,0)}{2.2}{bkAppleRed}}
  \foreach \x in {23,28,33} {\bkballoon{(\x,0)}{2.2}{bkRoute6}}
\end{tikzpicture}
\end{minipage}
\end{reidai}

\begin{jibun}
あかい きんぎょが 6ぴき、 くろい きんぎょが 2ひき。 あわせて なんびき？
\bkshikiran
\bkjans{6 ＋ 2 ＝ 8\quad 8ひき}
\end{jibun}

\begin{reidai}[ふえる]
こうえんで 5\R{人}{にん} あそんで います。 3\R{人}{にん} きました。 みんなで なん\R{人}{にん}？

\kangaekata 5\R{人}{にん}に 3\R{人}{にん} くわわった ぜんぶ。 たし\R{算}{ざん}。

\shiki 5 ＋ 3 ＝ 8 \kotaenext 8\R{人}{にん}
\end{reidai}

\begin{jibun}
どんぐりを 6こ ひろいました。 また 3こ ひろいました。 ぜんぶで なんこ？
\bkshikiran
\bkjans{6 ＋ 3 ＝ 9\quad 9こ}
\end{jibun}

\begin{reidai}[ものさしの うえを すすむ]
ものさしの 2の ところから 5つ すすむと、 どこ？
\begin{bkfigure}
\begin{tikzpicture}[x=1mm, y=1mm]
  \bkruler{(0,0)}{10.5}{10}
  \bkhop{(0,0)}{10.5}{2}{5}{bkBlue}
\end{tikzpicture}
\end{bkfigure}
\kangaekata \R{右}{みぎ}へ すすむのは たし\R{算}{ざん}。

\shiki 2 ＋ 5 ＝ 7 \kotaenext 7の ところ
\end{reidai}

\begin{jibun}
ものさしの 5の ところから 4つ すすむと、 どこ？
\bkshikiran[9\zw]
\bkjans{5 ＋ 4 ＝ 9\quad 9の ところ}
\end{jibun}

\begin{reidai}[10を こえる]
あめが 8こ。 5こ もらうと、 ぜんぶで なんこ？

\kangaekata 8は あと 2で 10。 5を 2と 3に わける。 8と 2で 10、 10と 3で 13。

\shiki 8 ＋ 5 ＝ 13 \kotaenext 13こ
\end{reidai}

\begin{jibun}
あかい いろがみが 9まい、 あおい いろがみが 3まい。 あわせて なんまい？\enspace{\small\color{bkOrange!58!black}ヒント：9は あと 1で 10}
\bkshikiran
\bkjans{9 ＋ 3 ＝ 12\quad 12まい}
\end{jibun}

%% ---------------- 第11回：10を こえる（もう1つ）・0を たす ----------------
\begin{reidai}[10を こえる（もう 1つ）]
クッキーが 7まい。 6まい やけました。 ぜんぶで なんまい？

\noindent\begin{minipage}[c]{0.56\linewidth}
\kangaekata 7は あと 3で 10。 6を 3と 3に わける。 7と 3で 10、 10と 3で 13。

\shiki 7 ＋ 6 ＝ 13 \kotaenext 13まい
\end{minipage}\hfill
\begin{minipage}[c]{0.42\linewidth}\centering
\begin{tikzpicture}[x=1mm, y=1mm]
  \bktenframe{(0,0)}{6}{7}{bkRoute2}{3}{bkRoute4}
  \foreach \x in {35,41,47} {\bkohajiki{(\x,3)}{2.3}{bkRoute4}}
  \node[anchor=north, font=\sffamily\bfseries\footnotesize] at (15,-1) {10};
  \node[anchor=north, font=\sffamily\bfseries\footnotesize] at (41,-1) {3};
\end{tikzpicture}
\end{minipage}
\end{reidai}

\begin{jibun}
ひよこが 6わ。 5わ うまれました。 ぜんぶで なんわ？\enspace{\small\color{bkOrange!58!black}ヒント：6は あと 4で 10}
\bkshikiran
\bkjans{6 ＋ 5 ＝ 11\quad 11わ}
\end{jibun}

\begin{reidai}[0を たす]
わなげを 2かい しました。 1かいめは 4てん、 2かいめは 0てん。 あわせて なんてん？

\kangaekata 2かいめは なにも ふえない。 0を たしても かずは かわらない。

\shiki 4 ＋ 0 ＝ 4 \kotaenext 4てん
\end{reidai}

\begin{jibun}
かごに たまごが 0こ。 5こ いれました。 たまごは なんこ？
\bkshikiran
\bkjans{0 ＋ 5 ＝ 5\quad 5こ}
\end{jibun}

\begin{jibun}[10を こえる（れんしゅう）]
あかい ビーズが 7こ、 あおい ビーズが 8こ。 あわせて なんこ？\enspace{\small\color{bkOrange!58!black}ヒント：7は あと 3で 10。 8を 3と 5に わけよう}
\bkshikiran
\bkjans{7 ＋ 8 ＝ 15\quad 15こ}
\end{jibun}



\clearpage
%% ---------------- p.52 まちがえやすいところ ----------------
\bksection{tsumazuki}{まちがえやすい ところ}

\begin{tsumazuki}{ふえたの？ へったの？}
すずめが 5わ います。 3わ とんで きました。 いま なんわ？\\
\bkbatsu\enspace へったと おもって 2わ\hspace{2\zw}\bkmaru\enspace 5 ＋ 3 ＝ 8\par
はじめの 5わに、 あとから 3わ くわわった。 「いま ぜんぶで なんわ？」と きいて いるから、 たし\R{算}{ざん}。
\end{tsumazuki}

\begin{tsumazuki}{とばす、 2かい かぞえる}
\begin{minipage}[c]{0.6\linewidth}
ばらばらの ものは、 しるしを つけながら かぞえよう。
\end{minipage}\hfill
\begin{minipage}[c]{0.36\linewidth}\centering
\begin{tikzpicture}[x=1mm, y=1mm]
  \foreach \x/\y/\c in {0/0/1,9/4/1,18/-1/1,4/-10/0,14/-9/0,24/-6/0,27/4/0} {
    \bkapple{(\x,\y)}{3.4}{bkAppleRed}
    \ifnum\c=1 \draw[bkBlue, line width=1.4pt, line cap=round] (\x-1.5,\y-0.5) -- (\x-0.2,\y-2) -- (\x+2.2,\y+1.5);\fi}
\end{tikzpicture}
\end{minipage}
\end{tsumazuki}

\begin{tsumazuki}{0を たすと 0？}
\bkbatsu\enspace 5 ＋ 0 ＝ 0\hspace{3\zw}\bkmaru\enspace 5 ＋ 0 ＝ 5\par
0は なにも ない かず。 なにも ふえないので、 5の ままです。
\end{tsumazuki}

\begin{tsumazuki}{10を わすれる}
\bkbatsu\enspace 8 ＋ 5 ＝ 3\hspace{3\zw}\bkmaru\enspace 8 ＋ 5 ＝ 13\par
8と 2で できた 「10の まとまり」を わすれないで。
\end{tsumazuki}



\clearpage
%% ---------------- p.53 つかわれているところ ----------------
\begin{butsuri}
たし\R{算}{ざん}は、 まちの なかで はたらいて います。
\begin{bkfigure}
\begin{tikzpicture}[x=1mm, y=1mm]
  % レジ
  \fill[white, rounded corners=3mm] (0,0) rectangle (59,36);
  \fill[bkGray!70, rounded corners=1.5mm] (12,11) rectangle (46,21);
  \fill[bkInk!85, rounded corners=1mm] (16,23) rectangle (42,32);
  \node[text=bkYellow, font=\sffamily\bfseries] at (29,27.5) {120 \R{円}{えん}};
  \foreach \x in {18,25,32,39} {\fill[white] (\x-2,14.5) rectangle ++(4,3);}
  \node[font=\sffamily\bfseries\small] at (29.5,5.5) {おみせの レジ};
  % えきの かいさつ（じどうかいさつき）
  \begin{scope}[shift={(63,0)}]
  \fill[white, rounded corners=3mm] (0,0) rectangle (59,36);
  \draw[bkInk, line width=1pt] (3,10.5) -- (56,10.5);
  \fill[bkGray!55, rounded corners=1mm] (6,10.5) rectangle (21,22) (38,10.5) rectangle (53,22);
  \fill[bkGray!85, rounded corners=0.6mm] (5,21) rectangle (22,23.2) (37,21) rectangle (54,23.2);
  \fill[bkRoute6!75, rounded corners=0.5mm] (15,23.2) rectangle (21,24.6) (38,23.2) rectangle (44,24.6);
  \fill[bkOrange, rounded corners=0.4mm] (21,15.5) rectangle (23.4,19.5) (35.6,15.5) rectangle (38,19.5);
  \bkperson{(29.5,10.5)}{17}{bkRoute2}
  \node[draw=bkInk, fill=white, rounded corners=1mm, font=\sffamily\bfseries\footnotesize, inner sep=1.5pt] at (29.5,31.5) {＋1 ＋1 ＋1};
  \node[font=\sffamily\bfseries\small] at (29.5,5.5) {えきの かいさつ};
  \end{scope}
\end{tikzpicture}
\end{bkfigure}
\begin{itemize}[label={\color{bkPurple}●}, leftmargin=1.4\zw, itemsep=0pt, topsep=0pt]
\item レジは、 かった ものの ねだんを ぜんぶ あわせて います。
\item えきの かいさつは、 きっぷや カードを とおすと 1\R{人}{にん} たして、 とおった \R{人}{ひと}の かずを かぞえます。
\end{itemize}
\end{butsuri}

\bkfukidashi{delta}{たし\R{算}{ざん}は、 この さきの \R{本}{ほん}でも ずっと つかうよ。 「たしあわせの \R{路線}{ろ|せん}」に のって いこう！}

\begin{bkfigure}
\begin{tikzpicture}[x=1mm, y=1mm]
  \draw[bkRoute2, line width=3pt] (10,0) -- (128,0);
  \node[anchor=west, inner sep=0pt] at (0,12) {\bkrouteicon[1.4em]{2}\hspace{0.3em}{\sffamily\bfseries\small\color{bkRoute2!80!black}たしあわせの \R{路線}{ろ|せん}}};
  \foreach \x/\c/\t/\s in {10/bkVol1/{1-4}/{たし\R{算}{ざん}},39.5/bkVol1/{1-7}/{\R{大}{おお}きな かず},69/bkVol1/{1-12}/{かけ\R{算}{ざん}},98.5/bkVol1/{1-21}/{ひろさ},128/bkVol2/{2-3}/{{\ltjsetparameter{jacharrange={+200}}＋－×÷}}} {
    \fill[white, draw=\c, line width=1.6pt] (\x,0) circle (2.6);
    \node[anchor=base, font=\sffamily\bfseries\footnotesize, text=\c!85!black] at (\x,-8.4) {\t};
    \node[anchor=base, font=\sffamily\footnotesize, text=bkInk] at (\x,-15.8) {\s};}
  \node[fill=bkTheme, text=white, rounded corners=1mm, font=\sffamily\bfseries\fontsize{8}{9}\selectfont, inner sep=1.5pt] at (10,5.5) {いま ここ};
\end{tikzpicture}
\end{bkfigure}

\begin{yattemiyou}[うちの なかで さがそう]
うちの なかで、 「あわせる」 ことや 「ふえる」 ことを さがそう。 （れい：げんかんの くつ）
\end{yattemiyou}

\clearpage
%% ---------------- p.54〜55 もんだい ----------------
\bkendproblems

\bkproblemgroup{きほん}{1}

\bkproblem けいさんを しましょう。
\par\vspace{2pt}
\noindent\hspace*{1\zw}\begin{tblr}{colspec={*{2}{Q[l,m,wd=62mm]}}, colsep=0pt, rowsep=3pt}
（1）2 ＋ 3 ＝ \bkkakikomi & （2）4 ＋ 5 ＝ \bkkakikomi \\
（3）6 ＋ 1 ＝ \bkkakikomi & （4）3 ＋ 7 ＝ \bkkakikomi \\
（5）9 ＋ 2 ＝ \bkkakikomi & （6）8 ＋ 6 ＝ \bkkakikomi \\
（7）7 ＋ 5 ＝ \bkkakikomi & （8）6 ＋ 6 ＝ \bkkakikomi \\
（9）5 ＋ 0 ＝ \bkkakikomi & （10）0 ＋ 4 ＝ \bkkakikomi \\
\end{tblr}
\bkpans{（1）5\quad（2）9\quad（3）7\quad（4）10\quad（5）11\quad（6）14\quad（7）12\quad（8）12\quad（9）5\quad（10）4}

\bkproblem あかい はなが 5ほん、 しろい はなが 3ぼん。 あわせて なんぼん？
\bkshikiran
\bkpans{5 ＋ 3 ＝ 8\quad 8ほん}

\bkproblem いけに かめが 4ひき います。 3びき きました。 ぜんぶで なんびき？
\bkshikiran
\bkpans{4 ＋ 3 ＝ 7\quad 7ひき}

\bkproblem えを みて、 たし\R{算}{ざん}の しきを かきましょう。
\begin{bkfigure}
\begin{tikzpicture}[x=1mm, y=1mm]
  \fill[bkTheme!6, rounded corners=3mm] (0,-8) rectangle (120,13);
  \foreach \x in {8,20,32,44,56,68} {\bkballoon{(\x,5)}{3}{bkAppleRed}}
  \foreach \x in {88,100,112} {\bkballoon{(\x,5)}{3}{bkRoute6}}
  \node[font=\sffamily\bfseries\small, text=bkInk] at (78,4) {と};
\end{tikzpicture}
\end{bkfigure}
\bkshikiran
\bkpans{6 ＋ 3 ＝ 9\quad 9こ}

\bkproblemgroup{チャレンジ}{1}

\bkproblem ものさしの 4の ところから 5つ すすむと、 どこに つく？
\begin{bkfigure}
\begin{tikzpicture}[x=1mm, y=1mm]
  \bkruler{(0,0)}{10.5}{10}
  \fill[bkTheme] (42,3.2) circle (1.6);
\end{tikzpicture}
\end{bkfigure}
\bkshikiran[9\zw]
\bkpans{4 ＋ 5 ＝ 9\quad 9の ところ}

\bkproblem バスに 7\R{人}{にん} のって います。 つぎの バスていで 6\R{人}{にん} のりました。 いま なん\R{人}{にん}？
\bkshikiran
\bkpans{7 ＋ 6 ＝ 13\quad 13\R{人}{にん}}


\begin{furikaeri}
\begin{checklist}
\item 「あわせる」「ふえる」を しきに かける。\hfill{\footnotesize\sffamily\color{bkTheme}→\pageref{bk:1-4:awaseru}ページ}
\item ものさしの うえで たし\R{算}{ざん}が できる。\hfill{\footnotesize\sffamily\color{bkTheme}→\pageref{bk:1-4:monosashi}ページ}
\item 10を こえる たし\R{算}{ざん}が できる。\hfill{\footnotesize\sffamily\color{bkTheme}→\pageref{bk:1-4:kuriagari}ページ}
\end{checklist}
チェックが つかなかったら、 →の ページに もどろう。
\end{furikaeri}

\bkfukidashi{sora}{ぜんぶ できたら、 おうちの \R{人}{ひと}に たし\R{算}{ざん}の もんだいを 1つ だして もらって、 といて みよう！}
\clearpage
